\chapter{任务二：结构化工具、执行合同与 Task Channel}
\label{chap:task2}

\section{约束}

结构化工具调用必须同时解决发现、参数、错误、安全和演进。一个只接受
\code{arg0/arg1} 的固定 switch 可以执行动作，却不能让模型可靠理解参数；一个只在
用户态校验 JSON Schema 的框架又无法阻止恶意 Guest 直接调用系统调用。我们因此把
模型友好的 rich description 放在 Guest，把紧凑、版本化、可强制的 typed ABI 放在
内核。

任务二只要求至少三项工具。当前内核目录包含 25 项工具，并提供 name/id 唯一映射、
V1 兼容描述、V2 typed descriptor 和结构化错误。数量不是主要贡献；关键是每条执行
路径具有清楚且不被混写的保证。

\section{工具发现和模型呈现分层}

\code{sys_tool_list} 返回定长 V2 descriptor，包括 version、size、tool id、name、schema、
flags 和 side-effect 类别。内核启动时验证目录完整性、id/name 唯一性和 schema 规则。
调用者可以用 name 或 id，但二者同时存在时必须指向同一工具；未知工具、未知参数、
缺失参数、类型错误、size/version 错误和缓冲区问题都返回稳定状态码。

Nexus 不把 25 项目录原样塞给模型。Guest runtime 从真实目录建立 role-filtered view，
补充 when-to-use、when-not-to-use、参数、结果、副作用和错误说明，再按查询类别返回。
运行时内部的 wait、Context 和 LLM 原语隐藏。这个分层使内核 schema 保持紧凑，模型
描述保持可读，同时不让 rich prompt 成为授权来源。

\section{Exploratory typed V2}

V2 请求携带显式 version、request size、tool name/id、parameter count 和 sized typed
parameter 数组。当前 scalar type 只有 \code{uint64} 与 string，降低教学内核解析面。
内核一次性复制参数，验证 key/type/数量后，才压平到共享执行器的
\code{arg0/arg1/64-byte payload}。压平是执行器内部兼容层，不改变 UAPI 的 typed
校验事实。

V2 适合自适应 Agent Loop：模型可以先查询文件，再根据候选数决定读取摘要或委派。
每次调用仍检查 role/capability、scope、schema、资源和工具自身不变量，但调用图没有
提前冻结。因此我们称它为 exploratory RPC，不把它写成 V3 相同的合同保证。

\section{Compact batch}

\code{agent_run()} 接受最多 64 个 compact \code{agent_op}，在一次 syscall 中按数组
顺序同步执行。它用于固定操作序列和 Context 压力验证，入口调用次数少，但不是 typed
KV batch、并行 batch 或 non-blocking runtime。任何一项结果仍写入对应 Context；
返回数组保留逐项状态。

这种兼容路径也是性能解释的必要边界。one-shot Task 实验比较的是等价 16 次空
\code{ECHO} 语义下的三种 transport sequence，而不是声称三条路径拥有相同 wire
format 或相同已知复制字节。

\section{ENFORCE V3 冻结执行合同}

高风险、可预先声明的工作流使用 V3。每个 lifecycle generation 最多冻结一份版本 1
合同，最多 24 个拓扑有序节点；predecessor 只能指向更小 node id。节点冻结 tool、
schema digest、required capability、input/output artifact type、deadline、attempt、
retry/cancel、provenance/effect manifest 和 exec/storage envelope。

V3 保留 V2 prefix，同时绑定 contract generation、node/attempt、source node、精确
predecessor Context sequence 和 32-byte fingerprint。ENFORCE 激活后，合同外调用和
受保护 direct syscall 在效果前被拒绝。每个 accepted \code{node+attempt} 使用稳定
完成槽；合法重试返回原 terminal，不重复副作用。每份合同最多 48 个 attempt 槽。

\section{Tool Phase Credit Lease}

合同节点还可锁定工具峰值资源。进入工具前，Phase Lease 从账户现有 U 中原子锁定
exec/storage envelope，状态按
\code{ADMITTED→ACTIVE→DEACTIVATED→SETTLED} 推进。分配器发布对象前取得 nonce claim；
失败对象把额度退回 locked U，成功对象继续持有同一 U，最终析构只结算一次。

Lease 不增加 quota，也不允许预测式超卖。它解决的是多类对象在一次工具内短时共同
达到峰值时，如何保证 envelope 在副作用前可用；未使用部分 terminal 时转为 F。
active phase 不能任意阻塞，全局 claim 和每 lifecycle phase 数都受合同 24 节点上限
约束。

\section{Typed Task Channel core}

Task Channel 借鉴 SQ/CQ 通信形状\cite{iouring}，但实现为 AgentOS 自有 ABI。每 Agent
按需映射一个 16-slot SQ 与一个 16-slot CQ，另有 request/resource 两个内核私有页，
共四页。SQE/CQE 均为 128 字节；内核先复制完整 SQE，再检查 ring/slot generation、
单调 request id、contract/node/attempt、schema、deadline、cancel link 和 typed handle，
避免共享页 TOCTOU。

私有 head/tail 是权威状态。CQ full 产生 backpressure；水位异常设置 sticky resync，
issuer 必须显式恢复。每个 accepted target 在 success/failure/cancel/timeout 中只发布
一个 terminal CQE；cancel 自己不产生第二个 CQE。completion 携带 Context sequence 与
Evidence ticket，但“one terminal CQE”不等于远程副作用的分布式 exactly-once。

当前 provider 同步、只接受 null input 和 output \code{NONE}；八槽 typed resource
表已实现 generation、owner、digest、producer 与 provenance 重验，但
\code{RESOURCE_IMPORT} fail closed，没有业务 payload/result backend。Nexus TASK 和模型
wire 都不走这条 channel。

\section{验证}

\code{agenttoolabi_ucore} 覆盖目录发现、name/id 冲突、未知工具、缺参、错参、边界
buffer 和 structured response。\code{agentcontract_ucore} 覆盖正常 DAG、deadline、
retry/cancel、source Context 与不可信 provenance 的 deterministic pre-effect denial。
\code{agenttask_ucore} 让 compact batch、scalar V3、SQ/CQ 分别执行 16 次空 ECHO，要求
同一 tool/status/Context-proof/evidence-proof/zero-result fingerprint，同时验证 full、
resync 和 retained-terminal cancel。

one-shot campaign 保留 96 条 16-op sequence。中位时延是 batch 561 us、scalar V3
2051 us、SQ/CQ 1620.5 us。它们是 sequence 级 wall-time 区间；逐 operation 的 Context
service-start interval 只有 10 ms tick 分辨率，不能均摊成伪造的单请求微秒。

\begin{evidencebox}[任务二结论]
我们交付了 25 项可发现工具和三类结构化执行机制。V2 提供自适应探索，V3 提供冻结
合同，Task Channel 提供有界 SQ/CQ core；当前 Task provider 和资源 backend 的限制在
接口、测试与文档中均 fail closed。
\end{evidencebox}
